% !TEX root = ../Thesis.tex

% =============================================================================
% KAPITEL 9 – Diskussion
% =============================================================================
 
\chapter{Diskussion}
 
Die Evaluation in Kapitel~8 hat die Ergebnisse der prototypischen Umsetzung an den Anforderungen aus Kapitel~3 gemessen und dabei sowohl funktionale Stärken als auch offene Lücken identifiziert. Das vorliegende Kapitel ordnet diese Befunde in den Kontext der wissenschaftlichen Literatur ein, bewertet die Übertragbarkeit auf vergleichbare Plattformen und grenzt den wissenschaftlichen vom praktischen Beitrag der Arbeit ab.
 
 
\section{Einordnung der Ergebnisse im Verhältnis zur Literatur}
 
Die Arbeit bestätigt mehrere Grundannahmen der Observability-Literatur auf konzeptioneller Ebene und konkretisiert sie zugleich um operative Erkenntnisse aus der prototypischen Umsetzung. Ein erheblicher Teil der herangezogenen Literatur behandelt Observability entweder auf konzeptioneller Ebene oder entlang konkreter Frameworks und Werkzeuge \parencite[S.~72034]{faseeha2025observability}, während die Integration heterogener Telemetriequellen und die Interoperabilität zwischen verschiedenen Werkzeugen als eigenständige Herausforderung beschrieben werden \parencite[S.~105]{sakinala2025observability}. Die vorliegende Arbeit ergänzt diese Perspektive um fallbezogene Befunde aus einer konkreten Implementierung.
 
Auf der Bestätigungsseite stehen drei Befunde. Erstens erweist sich die Kombination der drei Telemetriepfeiler Logs, Metriken und Traces als praktisch tragfähig. \textcite[S.~22]{albuquerque2025tracing} beschreiben \emph{Application Metrics} und \emph{Infrastructure Metrics} als komplementäre Entwurfsmuster, deren Zusammenspiel eine vollständige Sicht auf einen Dienst ermöglicht. Die vorliegende Arbeit stützt diesen Befund, indem sie zeigt, dass erst die Verknüpfung von Span-basierten Applikationsmetriken mit Exporter-basierten Infrastrukturmetriken in einem gemeinsamen Dashboard eine diagnostisch nutzbare Gesamtsicht erzeugt. Ohne diese Verknüpfung blieben die Infrastrukturindikatoren interpretationsarm, da sie zwar mögliche Engpässe anzeigen, nicht aber deren tatsächliche Ursache offenlegen \parencite[S.~20]{albuquerque2025tracing}.
 
Zweitens hat sich das Prinzip der symptomorientierten Alarmierung im Prototyp für den Latenz- und den Fehler-Alert bewährt, die jeweils nach einer Mehrsignal-Logik aufgebaut sind. Der Latenz-Alert verlangt neben der p95-Überschreitung gleichzeitig eine Mindest-Traffic-Rate als Vorbedingung und verhinderte in der Evaluation falsch-positive Auslösungen in Leerlaufphasen. Damit bestätigt die Arbeit die Empfehlung, Alarme auf nutzerseitig wahrnehmbare Zustände auszurichten und mehrere korrelierte Größen zu einem aussagekräftigeren Alarm zu kombinieren \parencite[S.~175]{vinnakota2025creating}. Einschränkend ist anzumerken, dass der Infrastruktur-Sättigungsalert in der prototypischen Regeldefinition vereinfacht umgesetzt wurde, sodass eine seiner Alternativbedingungen auch ohne gleichzeitige Servicewirkung feuern konnte. Die Mehrsignal-Logik ist dort konzeptionell angelegt, aber nicht durchgängig implementiert.
 
Drittens belegt die Arbeit, dass Auto-Instrumentierung über den OTel Java Agent einen niedrigschwelligen Einstieg in Distributed Tracing ermöglicht. \textcite[S.~52]{bissig2024unified} stellen fest, dass der mentale Aufwand für Auto-Instrumentierung in synthetischen Benchmarks gering ausfällt. Für frag.jetzt bestätigt sich dieser Befund insofern, als die Instrumentierung des Spring-Boot-Backends ohne Codeänderungen funktionierte und unmittelbar verwertbare Traces erzeugte.
 
Neben diesen Bestätigungen macht die Arbeit drei Aspekte sichtbar, die in der Literatur eher abstrakt behandelt werden und für den Fall frag.jetzt konkrete Gestalt annehmen. Der erste betrifft die operative Komplexität. Die Literatur beschreibt überwiegend konzeptionelle Modelle und empfohlene Vorgehensweisen. Die prototypische Umsetzung zeigt für den Fall frag.jetzt, dass infrastrukturelle und konfigurationsbezogene Herausforderungen einen erheblichen Einfluss auf die tatsächliche Nutzbarkeit eines Observability-Systems haben. Der OOM-Kill von Tempo, die Netzwerkprobleme in der Docker-Compose-Topologie und das Span-Rauschen durch die Auto-Instrumentierung traten als konkrete Problemkategorien auf, die in den herangezogenen Quellen nicht in dieser Spezifik behandelt werden. \textcite[S.~9]{Niedermaier2019onObservability} identifizieren zwar knappe Ressourcen und reaktive Implementierung als typische Herausforderungen (C7, C9), verbleiben aber auf einer übergeordneten Ebene. Die Kaskadeneffekte, die auf einem Single-Host-System durch den Observability-Stack selbst entstehen können, werden dort nicht adressiert.
 
Der zweite Aspekt betrifft die Signalqualität. Die Literatur diskutiert Sampling als Strategie zur Volumenreduktion \parencite[S.~10]{albuquerque2025tracing}. Die vorliegende Arbeit konkretisiert diesen Punkt, indem sie gezieltes Filtering auf Collector-Ebene einsetzt, das diagnostisch wertlose Spans verwirft, bevor sie gespeichert werden. Im Unterschied zum Sampling, das relevante Daten anteilig verlieren kann, erlaubt das Filtering eine inhaltlich begründete Selektion. Dieser Ansatz adressiert ein Problem, das \textcite[S.~52]{bissig2024unified} für reale Projekte unter dem Stichwort unerwarteten Verhaltens der Instrumentierung beschreiben, ohne eine Lösungsstrategie auf Collector-Ebene vorzuschlagen.
 
Der dritte Aspekt betrifft die Validierbarkeit von Alerting-Regeln. \textcite[S.~175--176]{vinnakota2025creating} fordern ausdrücklich, Alarme vor dem produktiven Einsatz rigoros zu testen, etwa durch synthetisches Monitoring, Fehlerinjektion und kontrollierte Störungssimulationen in Staging-Umgebungen. Die Evaluation zeigt exemplarisch, dass diese Forderung in der Praxis auf erhebliche Reibung stoßen kann. Der Fehler-Alert ließ sich definieren und konfigurieren, aber nicht zuverlässig zum Feuern bringen, weil die Interaktion zwischen dynamischem Abfragefenster, PromQL-Aggregationslogik und der Erfassungssemantik der OTel-Metriken ein eigenständiges technisches Problem darstellt. Die Literatur benennt die Notwendigkeit des Testens, beleuchtet jedoch die konkreten Schwierigkeiten, die sich aus dem Zusammenspiel von Metriksemantik, Query-Verhalten und Staging-Verifikation ergeben, nur begrenzt.
 
 
\section{Übertragbarkeit auf ähnliche cloud-native Plattformen}
 
Die Frage der Übertragbarkeit lässt sich in drei Stufen beantworten, die den Grad der Kontextabhängigkeit widerspiegeln.
 
\emph{Hoch übertragbar} sind drei methodische Elemente, die unabhängig von der spezifischen Architektur von frag.jetzt funktionieren. Die \emph{Signal-Matrix-Methodik} aus Kapitel~6.4, die für jeden Dienst systematisch die vier Golden Signals auf konkrete Messgrößen, Zwecke und Stakeholder-Rollen abbildet, ist als Strukturierungswerkzeug auf jedes verteilte System anwendbar. Sie setzt lediglich voraus, dass die beteiligten Dienste und ihre kritischen Interaktionen bekannt sind. Die \emph{Fünf-Felder-Runbook-Struktur} (Symptom, Kontext, Diagnose, Behebung, Eskalation) ist ebenfalls stackunabhängig und lässt sich auf beliebige Alerting-Systeme anwenden. Die \emph{symptomorientierte Alert-Logik}, die nutzerwirksame Zustände als Auslöser verwendet und interne Ursachenindikatoren als Debugging-Kontext reserviert, folgt einem Prinzip, das in der SRE-Literatur als stackübergreifend gültig beschrieben wird \parencite[Kap.~\emph{Symptoms Versus Causes}]{ewaschukMonitoringDistributedSystems}.
 
\emph{Bedingt übertragbar} ist der OTel-Agent-basierte Instrumentierungsansatz. Der OTel Java Agent funktioniert besonders reibungsarm in JVM-basierten Umgebungen. Für Anwendungen in anderen Laufzeitumgebungen, etwa Python oder Node.js, kann der Aufwand höher ausfallen und zusätzliche manuelle Anpassungen erfordern \parencite[S.~52]{bissig2024unified}. Das Grundprinzip, die Instrumentierung über einen herstellerneutralen Standard zu entkoppeln und dadurch die Bindung an ein bestimmtes Backend zu vermeiden, bleibt jedoch unabhängig von der Laufzeitumgebung gültig. \textcite[S.~72034--72035]{faseeha2025observability} bestätigen, dass die Interoperabilität zwischen heterogenen Plattformen eine der zentralen offenen Herausforderungen darstellt, die durch standardisierte Telemetrieformate wie OpenTelemetry adressiert wird.
 
\emph{Kontextabhängig} sind die konkreten Implementierungsentscheidungen. Der LGTM-Stack hat sich für das Single-Host-Szenario von frag.jetzt als geeignet erwiesen, stellt aber keine universelle Empfehlung dar. Für Organisationen mit Kubernetes-nativer Infrastruktur und höherem Telemetrievolumen kommen andere Stack-Kombinationen oder skalierbare Varianten infrage \parencite[S.~498]{gudimetla2025transitioning}. Die Docker-Compose-Topologie und die konkreten Schwellenwerte der Alerting-Regeln sind systemspezifisch. Schwellenwerte müssen aus den jeweiligen Lastprofilen und nichtfunktionalen Anforderungen des Zielsystems abgeleitet werden, wie bereits in Kapitel~4.2 als methodischer Grundsatz formuliert.
 
Zusammengefasst sind die methodischen Konzepte, also das \emph{Wie} der Ableitung, die strukturierende Matrix und die Alert-Logik, übertragbar. Die konkreten technischen Entscheidungen, also das \emph{Womit} und die spezifischen Grenzwerte, bleiben kontextgebunden.
 
 
\section{Wissenschaftlicher und praktischer Beitrag}
 
Die Arbeit entwickelt kein neues theoretisches Modell für Observability. Ihr Beitrag liegt in der methodisch fundierten Anwendung, Integration und Operationalisierung bestehender Konzepte für einen konkreten Anwendungskontext.
 
Der \emph{wissenschaftliche Beitrag} besteht in der methodisch transparenten Ableitungslogik, die von der Systemanalyse über die Anforderungsdefinition und die servicebezogene Operationalisierung der Golden Signals bis zur kriteriengestützten Evaluation reicht. Dieses Vorgehen verbindet mehrere in der Literatur häufig getrennt behandelte Konzepte zu einem integrierten Strategieentwurf. Dazu gehören die Golden Signals als Strukturierungsheuristik, das Design komplementärer Metrikebenen, die symptomorientierte Alarmierung und die rollenbasierte Informationsaufbereitung. Die explizite Rückbindung an Prüffragen und Erfüllungsgrade in der Evaluation stellt sicher, dass die Ergebnisse nachvollziehbar und an den Anforderungen messbar bleiben. Darüber hinaus liefert die Arbeit fallbezogene operative Hinweise auf Problemkategorien, etwa Ressourcenkonflikte des Observability-Stacks mit der beobachteten Anwendung, die in der vorwiegend konzeptionell argumentierenden Literatur bislang wenig Raum einnehmen.
 
Der \emph{praktische Beitrag} umfasst eine funktionsfähige Observability-Infrastruktur für frag.jetzt, die aus Instrumentierungskonfigurationen, Dashboard-Definitionen, Alerting-Regeln und Runbooks besteht. Die dokumentierten Herausforderungen bei Netzwerktopologie, Ressourcendimensionierung und Signalfilterung sind für Teams, die Observability erstmals in einer Docker-Compose-basierten Umgebung einführen, unmittelbar verwertbar. Die Signal-Matrix und die Runbook-Struktur stehen als wiederverwendbare Artefakte zur Verfügung, die auch ohne den spezifischen Technologie-Stack adaptiert werden können.
 
Beide Beitragsebenen stehen in einem komplementären Verhältnis: Die methodische Ableitung gibt dem Prototyp seine wissenschaftliche Nachvollziehbarkeit, und der Prototyp gibt den methodischen Konzepten ihre praktische Überprüfbarkeit.